\section{Síntesis y ampliación}

La idea central de esta sesión es: \textbf{las ideas básicas del RL clásico (gradiente de política, advantage estimation, la distinción entre on-policy y off-policy) también se aplican al razonamiento de los LLM y son fundamentales en él}.

Takeaways principales:
\begin{enumerate}
    \item El SFT está limitado por la escasez de datos; el ritmo de scaling que se obtiene solo con ampliar los datos es muy lento
    \item El RL, al aprovechar las muestras negativas y la señal de advantage, multiplica la eficiencia de datos por aproximadamente 8
    \item El RL en línea (como GRPO/PPO) evita el desplazamiento de distribución gracias al muestreo continuo
    \item Los modelos Thinking muestran que, con RL, pueden surgir conductas de razonamiento complejas a partir de una outcome reward simple
    \item La calidad del preentrenamiento del modelo base es una condición previa para que el RL tenga éxito
\end{enumerate}

\subsection{Tabla general de la evolución de los métodos}

\begin{table}[H]
\centering
\begin{tabular}{p{0.18\textwidth} p{0.2\textwidth} p{0.24\textwidth} p{0.22\textwidth}}
\toprule
Método & Datos/señal utilizados & Beneficio principal & Riesgo principal \\
\midrule
SFT / NTP & Problema + solución de referencia & Entrenamiento simple, estable y fácil de escalar & Limitado por el volumen de datos de alta calidad; tiende a aprender patrones superficiales \\
RFT & Muestras positivas on-policy + filtrado por resultado & Amplifica rápidamente las trayectorias correctas ya existentes & Los pasos espurios se refuerzan junto con los correctos; la generalización puede empeorar \\
Offline RL & Trayectorias erróneas + advantage / preference & Aprovecha las muestras negativas y mejora notablemente la eficiencia de datos & Depende de la política de rollout; hay desplazamiento de distribución \\
Online RL / GRPO / PPO & Muestreo de la política actual + outcome / dense reward & Optimiza directamente la distribución de conductas actual; adecuado para thinking model & Costo alto; sensible al diseño de la recompensa y a la calidad del modelo base \\
\bottomrule
\end{tabular}
\caption{Relación y compromisos entre los tipos de métodos de la Lecture 10}
\end{table}

\begin{knowledgebox}{Lecturas adicionales}
\begin{itemize}
    \item DeepSeek-R1: Incentivizing Reasoning Capability in LLMs via RL (2025)
    \item Kimi K1.5: Scaling Reinforcement Learning with LLMs (2025)
    \item Setlur et al., ``RL on Incorrect Synthetic Data Scales the Efficiency of LLM Math Reasoning by Eight-Fold,'' NeurIPS 2024
    \item Gandhi et al., ``Cognitive Behaviors that Enable Self-Improving Reasoners,'' arXiv 2025
    \item Setlur et al., ``Rewarding Progress: Scaling Automated Process Supervision for LLM Reasoning,'' arXiv 2024
\end{itemize}
\end{knowledgebox}

\end{document}
